\documentclass[11pt,a4paper]{article}
\usepackage[utf8]{inputenc}
\usepackage{amsmath,amsfonts,amssymb}
\usepackage{geometry}
\geometry{margin=1in}
\usepackage{hyperref}
\usepackage[many]{tcolorbox}
\usepackage{enumitem}
\usepackage{fancyhdr}

\pagestyle{fancy}
\fancyhf{}
\fancyfoot[L]{Universitas Bengkulu - Teknik Elektro}
\fancyfoot[R]{\thepage}
\def\footrule{{\color{gray}\hrule}}

\title{\vspace{-2cm}\textbf{Lembar Kerja Mahasiswa (Belajar Mandiri)\\Persamaan Diferensial - Minggu 11: Persamaan Gelombang 1D}}
\author{Novalio Daratha \& Muhammad Arfan\\Program Studi Teknik Elektro\\Universitas Bengkulu}
\date{2026}

\begin{document}

\maketitle

\begin{tcolorbox}[colback=blue!5!white,colframe=blue!75!black,title=Instruksi LKM]
Kerjakan secara berkelompok atau mandiri untuk memahami perambatan gelombang pada dawai dan media transmisi kelistrikan. Soal menguji dari level C1 hingga C6. Alokasikan waktu 45-60 menit.
\end{tcolorbox}

\section*{Soal-Soal Latihan}

\begin{enumerate}[label=\textbf{\arabic*.}]
    \item \textbf{(C1 - Mengingat)} Apa persamaan yang mengatur gelombang yang bergerak dengan kecepatan konstan $c$?
    \vspace{2.5cm}

    \item \textbf{(C2 - Memahami)} Bagaimana efek secara grafis fungsi $f(x-ct)$ saat waktu $t$ bertambah? Jelaskan pergeserannya.
    \vspace{3.5cm}

    \item \textbf{(C3 - Mengaplikasikan)} Terapkan kondisi batas $u(0,t)=0$ pada solusi pemisahan variabel $X(x)=C_1\cos(\lambda x)+C_2\sin(\lambda x)$. Tentukan parameter yang bernilai nol.
    \vspace{3.5cm}

    \item \textbf{(C4 - Menganalisis)} Hubungkan kecepatan rambat mekanik pada dawai ($c = \sqrt{T/\rho}$, di mana $T=$ tegangan, $\rho=$ massa jenis linier) dengan frekuensi vibrasi. Apa yang terjadi pada frekuensi jika dawai ditegangkan (T ditingkatkan)?
    \vspace{5cm}

    \item \textbf{(C5 - Mengevaluasi)} Suatu model gelombang memberikan kecepatan $c = 0$. Evaluasi keadaan fisik dari sistem ini, apakah ada propagasi atau pengiriman sinyal informasi pada kondisi tersebut?
    \vspace{5.5cm}

    \item \textbf{(C6 - Mencipta)} Ciptakan sebuah model sederhana untuk sinyal pulsa listrik berbentuk kotak tunggal yang berjalan di sepanjang saluran transmisi. Nyatakan fenomena transmisi ini secara matematis dengan solusi D'Alembert untuk pulsa sembarang $f(z)$!
    \vspace{6cm}
\end{enumerate}

\vfill
\begin{center}
\textbf{Semoga berhasil!}
\end{center}

\end{document}
